\clearpage
\scanpage{11}
\uuline{Problèmes-types à résoudre}

\reversemarginpar
\margin{\tikz[baseline]{\node (a) at (0,0) {\S\,7};\draw (0,0) ellipse (.38 and .59);}}
\fbox{a} \uline{Anneaux de définition} pour des \unsure{structures} sans \unsure{automorphisme|automorphismes}

\margin{\tikz[baseline]{\node at (0,0) {\S\S\,20,21};\draw (0,0) ellipse (.84 and .30);}}
b) \unsure{Hilbert}, Picard sans hypothèse projective
\normalmarginpar\margin{[\unsure{Résultats partiels}]}\reversemarginpar

\margin{\tikz[baseline]{\node at (0,0) {Chap \deleted{\unsure{V}}$\underline{\overline{\mathrm{VI}}}$};\draw (0,0) ellipse (.81 and .39);}}
\unsure{Dual} d'un schéma abélien non projectif
\normalmarginpar\margin{\fbox{OK. par Raynaud}}\reversemarginpar

\margin{\tikz[baseline]{\node at (0,0) {\S\S\,3,4};\draw (0,0) ellipse (.75 and .42);}}
\fbox{c} Passage au quotient sans hypothèse de platitude

\noindent\uuline{\hspace{3cm}}

% slanted
\margin{\tikz[baseline,remember picture]{\node[rotate=13] (mumford) at (0,0) {\unsure{Mumford}};\draw[rotate=13] (0,0) ellipse (.81 and .29);}}
\noindent\(\left[\begin{array}{@{}l@{}}
\tikz[baseline,remember picture]{\coordinate (itemd);}\text{d) Passage au quotient par }\unsure{\mathrm{Gl}\,n|\mathrm{Gl}_n|\mathrm{Gl}\,m}\text{, opérant ``sans point fixe''}\\[.5em]
\text{e) Modules de variétés plus ou moins quelconques.}\\
\qquad\tikz[baseline,remember picture]{\coordinate (abel);}\text{Cas des schémas abéliens}
\end{array}\right.\)
\tikz[remember picture,overlay]{\draw[gray,thin](mumford.east) -- (itemd); \draw[gray,thin](mumford.south east) to[out=-82,in=178] (abel);}

% slanted
\margin{\tikz[baseline]{\node[align=left,inner sep=1pt] at (0,0) {Mumford -\\\unsure{Seshadri} -\\\unsure{Narashiman|Narasimhan}};\draw (0,0) ellipse (1.01 and .67);}
pour \unsure{fibrés|fibres} vectoriels\thinspace; cas fini\thinspace: \unsure{rédiger|vérifier|vérifié}}
\noindent\(\left[\begin{array}{@{}l@{}}
\text{f) Classification de fibrés principaux }\inserted{\text{sur }X/Y}\ \text{(\unsure{par modules|par un autre})}\\
\quad\text{\unsure{à} groupe affine }G\text{ sur }Y\text{, on distingue les cas}\\
\qquad\tikz[baseline,remember picture]{\coordinate (abelg);}\,G\text{ commutatif}\\
\qquad G\text{ non néc. commutatif}
\end{array}\right.\)

\margin{\tikz[baseline,remember picture]{\node (seventeen) at (0,0) {\S\,17};\draw[rotate=14] (0,0) ellipse (.53 and .24);}}
\tikz[baseline,remember picture]{\coordinate (itemg);}g) Cas particulier\thinspace; \(G\) fini

\tikz[remember picture,overlay]{\draw[gray,thin] (seventeen.east) to[out=12,in=182] (abelg); \draw[gray,thin] (seventeen.east) to[out=-21,in=168] (itemg);}
% slanted
\margin{Pb mal posé\,?}
h) Cas particulier\thinspace: \(X=P^{1}_{Y},\quad G=\unsure{\mathrm{Gl}(n)|\mathrm{Gl}(m)}\).

\noindent\uuline{\hspace{3.7cm}}

\margin{Chap $\underline{\overline{\mathrm{VI}}}$\quad\tikz[baseline]{\node at (0,0) {oui};\draw[rotate=13] (0,0) ellipse (.53 and .35);}}
i) Les modules formels d'une variété abélienne \deleted{\unsure{$A$|$S$}} sont-ils \uline{simples} sur l'anneau de base\,?

\noindent\uuline{\hspace{3.8cm}}

\fbox{j)} ``Th. du carré'' dans \unsure{cas} non projectif \(\unsure{-\, -\, -\, -|\ldots}\)

% margin note written sideways
% Paper watermark through the centre.
\margin{EXTRA \unsure{P.J.A|P.J.} VELIN}
% archive sheet number: 10
